\begin{frame}{Inicialización greedy aleatorizada}
  \small
  \begin{columns}[T,onlytextwidth]
    \begin{column}{0.55\textwidth}
      \begin{block}{Construcción de planes iniciales}
        \begin{enumerate}
          \item Crear una plantilla con cursos pendientes y semestres sin asignar.
          \item Verificar condiciones necesarias y ventanas de prerrequisitos.
          \item Sortear un curso disponible y un semestre prometedor (RCL).
          \item Evaluar, reparar y conservar planes factibles distintos.
        \end{enumerate}
      \end{block}
    \end{column}
    \begin{column}{0.41\textwidth}
      \begin{exampleblock}{Feromona inicial de MMAS}
        \[M=\max_{p\in P_t}L_p(X^{\star})\]
        \[\tau_{\max}=\frac{1}{\rho M},\qquad
          \tau_{\min}=\gamma\,\tau_{\max}.\]
        \begin{itemize}
          \item $X^{\star}$: mejor plan factible inicial.
          \item $L_p$: créditos del semestre $p$; $M$: mayor carga.
          \item $\rho$: evaporación; $\gamma$: \texttt{tau-min-ratio}.
        \end{itemize}
        Todas las feromonas empiezan en $\tau_{\max}$.
      \end{exampleblock}
      % \begin{exampleblock}{Efecto esperado}
      %   La parte greedy guía la construcción; el azar aporta diversidad
      %   para iniciar la búsqueda MMAS.
      % \end{exampleblock}
      % \begin{alertblock}{Alternativas descartadas}
      %   Azar puro: muchas violaciones.\\[1mm]
      %   Greedy determinista: poca diversidad.\\[1mm]
      %   Método exacto: mayor costo para este hito.
      % \end{alertblock}
    \end{column}
  \end{columns}
  \vspace{0.2cm}
  {\footnotesize El mejor plan inicial se conserva como referencia; cada hormiga
  construye su propio plan desde un vector sin asignaciones.}
\end{frame}
